% Options for packages loaded elsewhere
\PassOptionsToPackage{unicode}{hyperref}
\PassOptionsToPackage{hyphens}{url}
\PassOptionsToPackage{dvipsnames,svgnames,x11names}{xcolor}
%
\documentclass[
  11pt,
  a4paper,
]{ltjsarticle}
\usepackage{amsmath,amssymb}
\usepackage{iftex}
\ifPDFTeX
  \usepackage[T1]{fontenc}
  \usepackage[utf8]{inputenc}
  \usepackage{textcomp} % provide euro and other symbols
\else % if luatex or xetex
  \usepackage{unicode-math} % this also loads fontspec
  \defaultfontfeatures{Scale=MatchLowercase}
  \defaultfontfeatures[\rmfamily]{Ligatures=TeX,Scale=1}
\fi
\usepackage{lmodern}
\ifPDFTeX\else
  % xetex/luatex font selection
\fi
% Use upquote if available, for straight quotes in verbatim environments
\IfFileExists{upquote.sty}{\usepackage{upquote}}{}
\IfFileExists{microtype.sty}{% use microtype if available
  \usepackage[]{microtype}
  \UseMicrotypeSet[protrusion]{basicmath} % disable protrusion for tt fonts
}{}
\makeatletter
\@ifundefined{KOMAClassName}{% if non-KOMA class
  \IfFileExists{parskip.sty}{%
    \usepackage{parskip}
  }{% else
    \setlength{\parindent}{0pt}
    \setlength{\parskip}{6pt plus 2pt minus 1pt}}
}{% if KOMA class
  \KOMAoptions{parskip=half}}
\makeatother
\usepackage{xcolor}
\usepackage[margin=24mm]{geometry}
\setlength{\emergencystretch}{3em} % prevent overfull lines
\providecommand{\tightlist}{%
  \setlength{\itemsep}{0pt}\setlength{\parskip}{0pt}}
\setcounter{secnumdepth}{5}
\ifLuaTeX
  \usepackage{selnolig}  % disable illegal ligatures
\fi
\IfFileExists{bookmark.sty}{\usepackage{bookmark}}{\usepackage{hyperref}}
\IfFileExists{xurl.sty}{\usepackage{xurl}}{} % add URL line breaks if available
\urlstyle{same}
\hypersetup{
  colorlinks=true,
  linkcolor={blue},
  filecolor={Maroon},
  citecolor={Blue},
  urlcolor={blue},
  pdfcreator={LaTeX via pandoc}}

\author{}
\date{}

\begin{document}

\hypertarget{chapter-02-ux884cux5217ux3068ux7ddaux5f62ux5199ux50cf}{%
\section{Chapter 02 ---
行列と線形写像}\label{chapter-02-ux884cux5217ux3068ux7ddaux5f62ux5199ux50cf}}

この章では、行列を表計算のような数表ではなく \textbf{空間を変形する関数}
として捉えます。

\hypertarget{ux884cux5217ux306eux4e09ux3064ux306eux9854}{%
\subsection{行列の三つの顔}\label{ux884cux5217ux306eux4e09ux3064ux306eux9854}}

\texttt{Ax} には同時に三つの意味があります。

\begin{enumerate}
\def\labelenumi{\arabic{enumi}.}
\tightlist
\item
  行ごとに見れば、各出力は入力との内積。
\item
  列ごとに見れば、出力は列ベクトルの一次結合。
\item
  全体として見れば、\texttt{R\^{}n→R\^{}m} の線形写像。
\end{enumerate}

この三つを切り替えられることが線形代数の強さです。

\hypertarget{ux6838ux3068ux50cfux304cux306aux305cux91cdux8981ux304b}{%
\subsection{核と像がなぜ重要か}\label{ux6838ux3068ux50cfux304cux306aux305cux91cdux8981ux304b}}

線形写像を理解する最小限の問いは

\begin{itemize}
\tightlist
\item
  どの入力が消えるか → \texttt{ker(A)}
\item
  どの出力に到達できるか → \texttt{Im(A)}
\end{itemize}

です。

この二つだけで、解の一意性、可逆性、情報損失をかなり説明できます。

例えば \texttt{z≠0} なのに \texttt{Az=0} なら

\[
A(x+z)=Ax
\]

です。つまり出力を見ても \texttt{x} と \texttt{x+z}
を区別できません。核は単なる方程式 \texttt{Ax=0}
の解集合ではなく、\textbf{観測で失われた自由度}です。

\hypertarget{rank-nullity-ux306eux8aadux307fux65b9}{%
\subsection{rank-nullity
の読み方}\label{rank-nullity-ux306eux8aadux307fux65b9}}

\[
\dim\ker A+\operatorname{rank}(A)=n
\]

を公式として暗記せず、入力の \texttt{n} 自由度が

\begin{itemize}
\tightlist
\item
  出力まで伝わる自由度
\item
  途中で消える自由度
\end{itemize}

へ分かれる保存則のように読みます。

\hypertarget{ux5230ux9054ux57faux6e96}{%
\subsection{到達基準}\label{ux5230ux9054ux57faux6e96}}

行列を見たとき「この行列は何を入力として何を出力する写像か」「何次元を保存し、何次元を失うか」を考えられることが目標です。

\end{document}
